%----------------------- Weekly Progress Document ------------------------
%
% Created by: Shane Reynolds 2020-02-29
%
%-------------------------------------------------------------------------

%-------------- Preamble
\documentclass[12pt]{article}
%!TEX root = weekly_progress_template.tex

\usepackage{geometry}
\geometry{a4paper, inner=2cm, outer=2cm, top=2cm, bottom=3cm}
  \pdfpagewidth=\paperwidth 
  \pdfpageheight=\paperheight
  % This acts as a failsafe to ensure things aren't stretched or moved when it's finally printed as a PDF.

\setlength{\parindent}{1em}  % Sets the length of the paragraph indent. Current setup has a an indent. Disable this if you activate the return line above.

% Double or one and a half spacing.
\usepackage{setspace}
  \onehalfspacing

\usepackage{booktabs}
% makes tables look sweet with top, mid, and bottom rules

\usepackage{graphicx}
  \DeclareGraphicsRule{.tif}{png}{.png}{`convert #1 `dirname #1`/`basename #1 .tif`.png}
% Graphics. Remove me and you won't have any figures, and that would be very boring.

\usepackage{fancyhdr}
% Want logos in your header? fancyhdr has got your back

%-------------- Include logo in header
\pagestyle{fancyplain} % use fancy for all pages except chapter start
\lhead{\includegraphics[height=5cm]{preamble/cdu_logo}} % left logo

%-------------- Function to create a small title block
\newcommand{\namelistlabel}[1]{\mbox{#1}\hfil}
\newenvironment{namelist}[1]{%1
\begin{list}{}
    {
        \let\makelabel\namelistlabel
        \settowidth{\labelwidth}{#1}
        \setlength{\leftmargin}{1.1\labelwidth}
    }
  }{%1
\end{list}} % this must be left as \input, \include wont work in the preamble

%-------------- Information For The Title Page
\title{	
		Thesis Progress Form\\
		CHARLES DARWIN UNIVERSITY\\
		College of Engineering, IT, and Environment
	  }
\author{}
\date{}

%------------------------ Main Document --------------------------
\begin{document}
	
	\maketitle
	
	\begin{namelist}{xxxxxxxxxxxx}
		\item[{\bf Name:}]
			Shane Reynolds
		\item[{\bf Unit:}]
			ENG720
		\item[{\bf Title:}]
			Automatic generation control of a two area power system using deep reinforcement learning
		\item[{\bf Supervisors:}]
			Charles Yeo \& Stefanija Klaric
		\item[{\bf Time \& Date:} 19/8/2020 \ @ 1pm]
			
	\end{namelist}
	
	\pagestyle{plain} % get rid of fancy headers
	\textheight = 565pt % hack to include page numbers
	
	%-------------- Sections
	\section{Progress since last meeting}
	\begin{itemize}
		\item Completed experiment in which agent memory was filled with tuned PID controller experience 50\% of the time. The controller learned a useful policy much quicker; however, was unable to 
		\item Developed stochastic demand profile to use in final experiments.
	\end{itemize}
	\section{Discussion Points}
	\begin{itemize}
		\item Outlined recent progress to CY, as described above.
		\item Outlined new experiments to CY, as detailed below.
	\end{itemize}
	\section{Plan until the next meeting}
	\begin{itemize}
		\item Run final experiments that will use a stochastic load demand profile to demonstrate that load frequency control can be achieved with a neural network using DDPG; however, this will not beat the current industry PID standard approach to control.
		\item Look at finding a way to present research with a narrative that shows training a neural network for load frequency control (and outperforming a PIC controller) is not possible. This would require some mathematical analysis.
	\end{itemize}
	%-------------- Supervisor sign-off
	\par
	\vspace{\fill}%
	\noindent\rule{0.4\linewidth}{0.5pt}%
	\vspace{1em}%
	\par
	\noindent\textbf{Supervisor}\vspace{1em}%
	\par
	\noindent\today

\end{document}